% ─────────────────────────────────────────────────────────────────────────────
% Thesis — Chapter 3: Implementation
% Sections 3.3 – 3.6
% Dan Herman, West University of Timișoara
% Supervisor: Lect. Dr. Adrian Spătaru
% ─────────────────────────────────────────────────────────────────────────────

% ─────────────────────────────────────────────────────────────────────────────
\section{Cowrie SSH Honeypot}
\label{sec:cowrie}
% ─────────────────────────────────────────────────────────────────────────────

The third stage of the attacker lifecycle is post-credential exploitation via
SSH\@. To capture this stage, the system deploys Cowrie~\cite{Son05}, a
medium-interaction SSH honeypot that emulates a full Unix shell environment
without exposing a real operating system.

\subsection{Deployment Architecture}

Cowrie runs directly on the host machine rather than inside a Docker container.
This design decision avoids a port-forwarding conflict: Docker's SSH daemon
occupies port~22, so placing Cowrie in a container would require a second level
of port mapping that complicates the bot-to-honeypot connection. Running bare on
the host allows Cowrie to bind directly to port~2222 and be reached by all four
bot containers via the host gateway address (\texttt{host.docker.internal}).

The honeypot is configured in \emph{shell} backend mode, meaning all commands
are handled by Cowrie's built-in command emulator rather than forwarded to a
real system. This is appropriate for the thesis scenario because the goal is
behavioural logging, not high-fidelity command execution. The relevant
\texttt{cowrie.cfg} parameters are shown in Table~\ref{tab:cowrie_cfg}.

\begin{table}[h]
\centering
\caption{Key Cowrie configuration parameters.}
\label{tab:cowrie_cfg}
\begin{tabular}{ll}
\hline
\textbf{Parameter} & \textbf{Value} \\
\hline
\texttt{hostname}              & \texttt{CORP-SRV-01} \\
\texttt{backend}               & \texttt{shell} \\
\texttt{auth\_class}           & \texttt{UserDB} \\
\texttt{listen port}           & \texttt{2222} \\
\texttt{idle\_timeout}         & 180 s \\
\texttt{authentication\_timeout} & 120 s \\
\texttt{ttylog}                & enabled \\
\texttt{logtype}               & rotating JSON \\
\hline
\end{tabular}
\end{table}

The hostname \texttt{CORP-SRV-01} is chosen to reinforce the corporate
deception established by the phishing page. An attacker who steals credentials
from what appears to be a corporate login portal and subsequently reaches a
server named \texttt{CORP-SRV-01} is more likely to proceed with reconnaissance
than to abort, increasing the volume and quality of captured behaviour.

\subsection{Credential Configuration}

Authentication is controlled by \texttt{userdb.txt}, which lists every
username–password pair that Cowrie will accept. The file contains nine entries:
the five planted corporate identities that are also seeded in the Flask decoy
and PostgreSQL (\texttt{john.mitchell}, \texttt{sarah.chen}, \texttt{admin},
\texttt{it.support}, \texttt{j.smith}), two common default credentials
(\texttt{root/123456}, \texttt{admin/admin}), and a wildcard entry
(\texttt{*:0:*}) that accepts any remaining credential pair. The wildcard is
important for realism: it ensures that scanners probing with arbitrary
credentials are also admitted, generating additional SSH session data rather
than being silently rejected at the authentication stage.

\subsection{Fake Filesystem}

Cowrie's \texttt{honeyfs/} directory overlays a set of files onto the emulated
filesystem that are designed to elicit further attacker interaction. Three files
are of particular interest:

\begin{itemize}
    \item \textbf{\texttt{/etc/passwd}} — Contains eleven entries including the
    five planted corporate users with realistic home directories and login
    shells. This file is the first target of account-discovery commands such as
    \texttt{cat /etc/passwd} and \texttt{awk '\{print \$1\}' /etc/passwd}.

    \item \textbf{\texttt{/etc/secrets.txt}} — A simulated internal
    configuration file containing fictitious database credentials, AWS access
    keys, and SMTP passwords. Planting credentials of this kind follows the
    deception principle of \emph{breadcrumbing}~\cite{LD08}: each secret gives
    the attacker a reason to continue and to attempt lateral movement, extending
    the session duration and command count captured by the logger.

    \item \textbf{\texttt{/home/john.mitchell/.bash\_history}} — A pre-populated
    command history that includes \texttt{ssh admin@10.0.0.15},
    \texttt{mysql -h 10.0.0.15}, and \texttt{scp} transfers to an internal
    backup server. These entries simulate the artefacts left by a legitimate
    administrator and encourage sophisticated attackers to investigate internal
    network topology.
\end{itemize}

\subsection{Logging}

Cowrie produces two log streams: a human-readable text log and a structured
JSON log (\texttt{var/log/cowrie/cowrie.json}). The JSON format is used
exclusively in this system because it is directly ingestible by Filebeat without
a parsing step. Each JSON event records the event type (\texttt{cowrie.login.success},
\texttt{cowrie.command.input}, \texttt{cowrie.session.closed}, etc.), the source
IP, session identifier, timestamp, and — for command events — the exact string
typed by the attacker. TTY logging is also enabled, producing a full terminal
replay of every session in \texttt{var/lib/cowrie/tty/}.

The \texttt{session\_id} field written into SSH sessions by the attacker bot is
the same UUID generated during the Flask phishing phase. This linkage is the
core contribution of the system and is discussed further in
Section~\ref{sec:ml_design}.


% ─────────────────────────────────────────────────────────────────────────────
\section{Attacker Bot Simulation}
\label{sec:bots}
% ─────────────────────────────────────────────────────────────────────────────

Evaluating a honeypot in a controlled laboratory environment requires a source
of realistic attacker traffic. Four automated bot profiles are implemented in
\texttt{bot\_manager.py}, each parameterised from real-world data and
running in its own isolated Docker container and network subnet.

\subsection{Justification for Synthetic Traffic}

The use of synthetic traffic in honeypot research is established practice when
live exposure is not feasible~\cite{BCF12}. The behavioural parameters of each
profile are derived from two authoritative sources:

\begin{enumerate}
    \item \textbf{CyberLab Honeynet Dataset}~\cite{RF06} (Zenodo record
    3687527) — 293 real SSH honeypot sessions from which login timing
    distributions, session duration medians, and login attempt counts are
    extracted.
    \item \textbf{MITRE ATT\&CK} — Techniques T1078 (Valid Accounts), T1059
    (Command and Scripting Interpreter), T1087 (Account Discovery), and T1005
    (Data from Local System) are used to select the command sequences executed
    by each profile.
\end{enumerate}

This dual grounding means that while the traffic is synthetic, the distributions
it samples from are empirically validated, satisfying the reproducibility
requirement of the experimental design.

\subsection{Two-Phase Attack Model}

Each bot cycle executes two phases that mirror the complete attacker lifecycle:

\begin{description}
    \item[Phase 1 — Victim bot] Simulates a user who receives a phishing email,
    visits the Flask decoy, submits credentials, and completes MFA\@. The bot
    reads the MFA code from the MailHog API, replicating the account takeover
    step without any external network dependency.

    \item[Phase 2 — Attacker bot] After a profile-specific delay, reads the
    stolen credentials from PostgreSQL and opens an SSH connection to Cowrie on
    port~2222. It then executes the command sequence defined for its profile.
\end{description}

Both phases share the same \texttt{session\_id} (a UUID generated at the start
of Phase~1). This identifier is written into \texttt{phishing\_sessions},
\texttt{ssh\_sessions}, and \texttt{ml\_features}, enabling every event in the
attacker lifecycle to be traced to a single origin.

\subsection{Profile Definitions}

The four profiles and their key parameters are summarised in
Table~\ref{tab:bot_profiles}. Full profile definitions are given below.

\begin{table}[h]
\centering
\caption{Attacker bot profile parameters derived from CyberLab distributions and MITRE ATT\&CK.}
\label{tab:bot_profiles}
\begin{tabular}{lllll}
\hline
\textbf{Profile} & \textbf{Login delay} & \textbf{MFA delay} & \textbf{Commands} & \textbf{MITRE technique} \\
\hline
\texttt{scanner}            & 0.1–0.5 s  & 0.1–0.3 s  & 4  & T1087 \\
\texttt{credential\_stuffer} & 0.5–2.0 s  & 0.2–0.5 s  & 4  & T1110 \\
\texttt{human\_like}         & 3.0–12.0 s & 8.0–30.0 s & 11 & T1059, T1005 \\
\texttt{sophisticated}       & 8.0–25.0 s & 15.0–45.0 s & 16 & T1078, T1087, T1005 \\
\hline
\end{tabular}
\end{table}

\paragraph{Scanner.}
Parameterised from 34 CyberLab sessions with a median duration of 25.4~s and
zero login attempts. The scanner submits all five credential pairs in rapid
succession (0.1–0.5~s per field), uses non-browser user agents
(\texttt{python-requests}, \texttt{curl}, \texttt{Nikto}), and executes four
basic account-discovery commands (\texttt{whoami}, \texttt{id},
\texttt{uname -a}, \texttt{cat /etc/passwd}) with inter-command delays of
0.1–0.5~s. This profile represents automated vulnerability scanning rather than
targeted intrusion.

\paragraph{Credential stuffer.}
Parameterised from 4 CyberLab sessions with a mean of 23.75 login attempts.
It submits ten credential pairs at 0.5–2.0~s per field, consistent with
list-based brute-force tooling. The command set targets privilege escalation
indicators (\texttt{cat /etc/shadow}, \texttt{ls /home}).

\paragraph{Human-like.}
Parameterised from 148 CyberLab sessions, the largest cluster, with a median
session duration of 52.1~s. Login delays of 3–12~s and MFA delays of 8–30~s
reflect the timing of a human operator reading and typing. The command sequence
of eleven entries includes network enumeration (\texttt{netstat -tulpn},
\texttt{ifconfig}) and data collection (\texttt{cat /etc/secrets.txt}),
consistent with MITRE ATT\&CK T1059 and T1005\@.

\paragraph{Sophisticated.}
Represents an Advanced Persistent Threat (APT) actor. Login delays of 8–25~s
and MFA delays of 15–45~s indicate deliberate, unhurried operation. The
sixteen-command sequence includes persistent access investigation
(\texttt{crontab -l}, \texttt{cat /etc/crontab}), broad configuration file
discovery (\texttt{find / -name '*.conf'}), and lateral movement indicators
(\texttt{ss -tulpn}). The profile executes a single login attempt and uses only
a legitimate browser user agent, minimising the automated footprint visible to
anomaly detectors.

\subsection{Network Isolation}

Each bot runs in a dedicated Docker subnet (Table~\ref{tab:bot_subnets}),
ensuring that the source IP address recorded in every session genuinely
differs per profile. This is essential for ML feature validity: if all bots
shared a subnet, the classifier could trivially memorise IP prefixes rather than
learning from behavioural features.

\begin{table}[h]
\centering
\caption{Docker subnet assignments for bot isolation.}
\label{tab:bot_subnets}
\begin{tabular}{lll}
\hline
\textbf{Profile} & \textbf{Subnet} & \textbf{Bot IP} \\
\hline
\texttt{scanner}            & \texttt{172.20.0.0/24} & \texttt{172.20.0.5} \\
\texttt{credential\_stuffer} & \texttt{172.21.0.0/24} & \texttt{172.21.0.5} \\
\texttt{human\_like}         & \texttt{172.22.0.0/24} & \texttt{172.22.0.5} \\
\texttt{sophisticated}       & \texttt{172.23.0.0/24} & \texttt{172.23.0.5} \\
\hline
\end{tabular}
\end{table}


% ─────────────────────────────────────────────────────────────────────────────
\section{Log Collection and Visualisation}
\label{sec:elk}
% ─────────────────────────────────────────────────────────────────────────────

Centralised log aggregation is implemented using the Elastic
Stack~\cite{Hol06}: Filebeat ships logs from both honeypot components to
Elasticsearch, where they are indexed and made available for analysis in Kibana.

\subsection{Filebeat Configuration}

Filebeat is deployed as a Docker service with bind-mount access to both log
directories. It is configured with two inputs and a single Elasticsearch output,
as shown in Table~\ref{tab:filebeat}.

\begin{table}[h]
\centering
\caption{Filebeat input and index configuration.}
\label{tab:filebeat}
\begin{tabular}{lll}
\hline
\textbf{Source} & \textbf{Log path} & \textbf{Elasticsearch index} \\
\hline
Flask phishing decoy & \texttt{/logs/flask/*.json}        & \texttt{honeypot-flask-\{date\}} \\
Cowrie SSH honeypot  & \texttt{/logs/cowrie/cowrie.json}  & \texttt{honeypot-cowrie-\{date\}} \\
\hline
\end{tabular}
\end{table}

Both inputs use \texttt{json.keys\_under\_root: true}, which promotes all JSON
fields to the top level of the Elasticsearch document rather than nesting them
under a \texttt{json} key. This simplifies Kibana query construction. Each
document is tagged with a \texttt{log\_source} field (\texttt{flask} or
\texttt{cowrie}) and a \texttt{honeypot\_component} field
(\texttt{phishing\_decoy} or \texttt{ssh\_honeypot}), enabling index-level
routing and per-component filtering within a single Kibana data view.

Index lifecycle management is disabled (\texttt{setup.ilm.enabled: false})
because the dataset is bounded in size; a custom index template matching
\texttt{honeypot-*} is applied instead to ensure consistent field mappings
across both index families.

\subsection{Elasticsearch}

Elasticsearch~8.13.0 runs as a single-node cluster on port~9200 within the
\texttt{honeypot-net} Docker network. The single-node configuration is
sufficient for the thesis scale and eliminates the coordination overhead of a
multi-node deployment. Health is confirmed via the cluster health API\@:

\begin{verbatim}
GET http://localhost:9200/_cluster/health
{"status":"green","number_of_nodes":1,...}
\end{verbatim}

\subsection{Kibana}

Kibana~8.13.0 is accessible on port~5601 and provides real-time dashboards over
the collected honeypot data. A single data view covering the \texttt{honeypot-*}
index pattern is configured, allowing analysts to query both Flask and Cowrie
events in the same interface and correlate them by \texttt{session\_id} or
timestamp. Key investigative queries include filtering by
\texttt{log\_source: cowrie} to isolate SSH sessions and by
\texttt{event.type: cowrie.command.input} to extract the full command history
for any individual session.


% ─────────────────────────────────────────────────────────────────────────────
\section{Machine Learning Pipeline}
\label{sec:ml_design}
% ─────────────────────────────────────────────────────────────────────────────

The ML pipeline is the analytical centrepiece of the system. Its purpose is to
classify each recorded session along two independent axes simultaneously,
constituting a multi-label classification problem~\cite{MDR25}.

\subsection{Problem Formulation}

Each session in \texttt{ml\_features} is assigned two ground-truth labels set
by the bot manager at collection time:

\begin{itemize}
    \item \textbf{\texttt{label\_bot\_type}} $\in$ \{\texttt{scanner},
    \texttt{credential\_stuffer}, \texttt{human\_like},
    \texttt{sophisticated}\} — identifies \emph{who} the attacker is.
    \item \textbf{\texttt{label\_attack\_stage}} $\in$
    \{\texttt{cred\_submission}, \texttt{ssh\_attack}\} — identifies
    \emph{what phase} of the lifecycle the session belongs to.
\end{itemize}

The two labels are treated independently: one Random Forest classifier is
trained per label, both using the same feature vector. This architecture allows
the classifiers to be evaluated separately and makes it possible to deploy them
in any combination without retraining.

\subsection{Feature Engineering}

Six features are extracted for each session, spanning both the phishing and SSH
stages of the lifecycle (Table~\ref{tab:features}). Two columns present in the
database schema (\texttt{user\_agent\_entropy} and
\texttt{inter\_cmd\_delay\_std\_s}) were found to be entirely unpopulated in the
collected dataset and are excluded from training. Sessions with the label
\texttt{phishing\_visit} are also excluded because no sessions of this type were
generated; the two observed stages are \texttt{cred\_submission} and
\texttt{ssh\_attack}.

\begin{table}[h]
\centering
\caption{ML feature set used for both classifiers.}
\label{tab:features}
\begin{tabular}{lll}
\hline
\textbf{Feature} & \textbf{Source stage} & \textbf{Description} \\
\hline
\texttt{time\_on\_page\_ms}       & Phishing & Total time from page load to form submit \\
\texttt{typing\_delay\_ms}        & Phishing & Mean delay between keystrokes in form fields \\
\texttt{field\_corrections}       & Phishing & Number of backspace/correction events \\
\texttt{command\_count}           & SSH      & Total commands executed in session \\
\texttt{session\_duration\_s}     & SSH      & Total SSH session duration in seconds \\
\texttt{inter\_cmd\_delay\_avg\_s} & SSH     & Mean delay between consecutive commands \\
\hline
\end{tabular}
\end{table}

For \texttt{cred\_submission} rows the SSH features are set to zero; for
\texttt{ssh\_attack} rows the phishing features are set to zero. This encoding
is deliberate: the zero values are themselves a discriminating signal — a session
with \texttt{time\_on\_page\_ms = 0} cannot be a phishing session, which is
information the classifier can exploit.

\subsection{Classifier Design}

Both classifiers use the \texttt{RandomForestClassifier} from
\texttt{scikit-learn}~\cite{MDR25} with the following configuration:

\begin{itemize}
    \item 200 estimators with no depth limit (\texttt{max\_depth=None}).
    \item \texttt{class\_weight='balanced'} — sample weights are inversely
    proportional to class frequency, compensating for the imbalance between the
    large \texttt{scanner}/\texttt{credential\_stuffer} classes and the smaller
    \texttt{human\_like}/\texttt{sophisticated} classes without discarding data.
    \item \texttt{random\_state=42} for reproducibility.
\end{itemize}

\subsection{Evaluation Protocol}

Model performance is estimated using 5-fold stratified cross-validation, which
preserves the class distribution in each fold. Three metrics are reported:
accuracy, weighted F1 (accounts for class imbalance by weighting each class by
its support), and macro F1 (treats all classes equally, penalising poor
performance on minority classes). The macro F1 is the primary evaluation metric
because it highlights classifier behaviour on the underrepresented
\texttt{human\_like} and \texttt{sophisticated} profiles, which are the most
forensically interesting cases.

\subsection{Baseline Comparison}

To contextualise the results, performance is compared against published
classification results on the CICIDS2017 dataset~\cite{BCF12}, the standard
benchmark for network intrusion detection. CICIDS2017 contains approximately
500\,000 labelled network flows across 14 attack categories and is routinely
used to evaluate Random Forest classifiers for intrusion detection. Reported
macro F1 scores on CICIDS2017 for Random Forest range from 0.89 to 0.97
depending on the feature set and preprocessing applied. The honeypot dataset
presents a harder problem in one respect — four-class behavioural profiling
rather than binary anomaly detection — and an easier one in another — ground
truth is exact because labels are set by the generating bots rather than
inferred by analysts.
